\documentclass[12pt]{article}
\usepackage[a4paper,margin=1in]{geometry}
\usepackage[utf8]{inputenc}
\usepackage[T1]{fontenc}
\usepackage{amsmath,amssymb}
\usepackage{textcomp}
\usepackage{parskip}
\begin{document}
\begin{center}
{\Large \textbf{1999 STEP I}}
\end{center}

\subsection*{Page 1}
\begin{flushleft}
STEP I, 1999 2\par
\bigskip\textbf{Section A: Pure Mathematics}\par
\bigskip\textbf{Question 1.} How many integers greater than or equal to zero and less than a million are not divisible by\par
2 or 5? What is the average value of these integers?\par
How many integers greater than or equal to zero and less than 4179 are not divisible by 3\par
or 7? What is the average value of these integers?\par
\bigskip\textbf{Question 2.} A point moves in the x-y plane so that the sum of the squares of its distances from the three\par
fixed points (x1, y1), (x2, y2), and (x3, y3) is always a . Find the equation of the locus of\par
the point and interpret it geometrically. Explain why a cannot be less than the sum of the\par
squares of the distances of the three points from their centroid.\par
\[[The centroid has coordinates (x¯, y¯) where 3x¯ = x^1 + x^2 + x^3, 3y¯ = y^1 + y^2 + y^3. ]\]
\bigskip\textbf{Question 3.} The n positive numbers x1, x2, . . . , x\par
n\par
\[, where n > 3, satisfy\]
1 1 1\par
\[x^1 = 1 + , x^2 = 1 + , . . . , x n-1 = 1 + ,\]
\[x^2 x^3 x\]
n\par
and also\par
\[x = 1 + .\]
n\par
\[x^1\]
Show that\par
\textbf{(i)} x1, x2, . . . , x\par
n\par
\[> 1,\]
\[x^2 - x^3\]
\textbf{(ii)} x1 - x2 = - ,\par
\[x2x3\]
\textbf{(iii)} x1 = x2 = \cdots = x\par
n\par
.\par
\[Hence find the value of x^1.\]
\bigskip\textbf{Question 4.} Sketch the following subsets of the x-y plane:\par
\textbf{(i)} |x| + |y| 6 1 ;\par
\textbf{(ii)} |x - 1| + |y - 1| 6 1 ;\par
\textbf{(iii)} |x - 1| - |y + 1| 6 1 ;\par
\textbf{(iv)} |x| |y - 2| 6 1 .\par
\end{flushleft}
\subsection*{Page 2}
\begin{flushleft}
STEP I, 1999 3\par
\bigskip\textbf{Question 5.} For this question, you may use the following approximations, valid if θ is small: sin θ \approx θ and\par
cos θ \approx 1 - θ /2 .\par
A satellite X is directly above the point Y on the Earth's surface and can just be seen (on\par
the horizon) from another point Z on the Earth's surface. The radius of the Earth is R and\par
the height of the satellite above the Earth is h.\par
\textbf{(i)} Find the distance d of Z from Y along the Earth's surface.\par
\textbf{(ii)} If the satellite is in low orbit (so that h is small compared with R), show that\par
d \approx k(Rh)\par
\[1/2\]
,\par
where k is to be found.\par
\textbf{(iii)} If the satellite is very distant from the Earth (so that R is small compared with h),\par
show that\par
d \approx aR + b(R /h),\par
where a and b are to be found.\par
\bigskip\textbf{Question 6.} (i) Find the greatest and least values of bx + a for -10 6 x 6 10, distinguishing carefully\par
\[between the cases b > 0, b = 0 and b < 0.\]
\textbf{(ii)} Find the greatest and least values of cx 2 + bx + a, where c > 0, for -10 6 x 6 10,\par
distinguishing carefully between the cases that can arise for different values of b and c.\par
\bigskip\textbf{Question 7.} Show that sin(k sin -1 x), where k is a constant, satisfies the differential equation\par
d y dy\par
2 2\par
\[(1 - x ) - x + k y = 0. (∗)\]
\[dx2 dx\]
\[In the particular case when k = 3, find the solution of equation (∗) of the form\]
3 2\par
\[y = Ax + Bx + Cx + D,\]
dy\par
\[that satisfies y = 0 and = 3 at x = 0.\]
dx\par
Use this result to express sin 3θ in terms of powers of sin θ.\par
\end{flushleft}
\subsection*{Page 3}
\begin{flushleft}
STEP I, 1999 4\par
\bigskip\textbf{Question 8.} The function f satisfies 0 6 f (t) 6 K when 0 6 t 6 x. Explain by means of a sketch, or\par
otherwise, why\par
x\par
0 6 f (t) dt 6 Kx.\par
\[\int\]
t\par
\[By considering dt, or otherwise, show that, if n > 1,\]
0 n(n - t)\par
\[\int\]
n 1 1 1\par
0 6 ln - 6 -\par
n - 1 n n - 1 n\par
µ ¶\par
and deduce that\par
N\par
0 6 ln N - 6 1.\par
n\par
\[n=2\]
X\par
Deduce that as N \to \infty\par
N\par
\to \infty.\par
n\par
\[n=1\]
X\par
10 30\par
\[Noting that 2 = 1024, show also that if N < 10 then\]
N\par
\[< 101.\]
n\par
\[n=1\]
X\par
\end{flushleft}
\subsection*{Page 4}
\begin{flushleft}
STEP I, 1999 5\par
\bigskip\textbf{Section B: Mechanics}\par
\bigskip\textbf{Question 9.} A tortoise and a hare have a race to the vegetable patch, a distance X kilometres from the\par
starting post, and back. The tortoise sets off immediately, at a steady v kilometers per hour.\par
The hare goes to sleep for half an hour and then sets off at a steady speed V kilometres per\par
hour. The hare overtakes the tortoise half a kilometre from the starting post, and continues\par
on to the vegetable patch, where she has another half an hour's sleep before setting off for the\par
return journey at her previous pace. One and quarter kilometres from the vegetable patch,\par
she passes the tortoise, still plodding gallantly and steadily towards the vegetable patch. Show\par
that\par
\[V =\]
4X - 9\par
and find v in terms of X.\par
Find X if the hare arrives back at the starting post one and a half hours after the start of the\par
race.\par
\bigskip\textbf{Question 10.} A particle is attached to a point P of an unstretched light uniform spring AB of modulus of\par
elasticity λ in such a way that AP has length a and P B has length b. The ends A and B of\par
the spring are now fixed to points in a vertical line a distance l apart, The particle oscillates\par
along this line. Show that the motion is simple harmonic. Show also that the period is the\par
same whatever the value of l and whichever end of the string is uppermost.\par
\bigskip\textbf{Question 11.} The force of attraction between two stars of masses m1 and m2 a distance r apart is γm1m2/r 2 .\par
The Starmakers of Kryton place three stars of equal mass m at the corners of an equilateral\par
triangle of side a. Show that it is possible for each star to revolve round the centre of mass of\par
the system with angular velocity (3γm/a\par
)\par
\[1/2\]
.\par
Find a corresponding result if the Starmakers place a fourth star, of mass λm, at the centre\par
of mass of the system.\par
\end{flushleft}
\subsection*{Page 5}
\begin{flushleft}
STEP I, 1999 6\par
\bigskip\textbf{Section C: Probability and Statistics}\par
\bigskip\textbf{Question 12.} (i) Prove that if x > 0 then x + x -1 > 2.\par
I have a pair of six-faced dice, each with faces numbered from 1 to 6. The probability\par
of throwing i with the first die is q and the probability of throwing j with the second\par
i\par
die is r (1 6 i, j 6 6). The two dice are thrown independently and the sum noted. By\par
j\par
considering the probabilities of throwing 2, 12 and 7, show the sums 2, 3, . . . , 12 are\par
not equally likely.\par
\textbf{(ii)} The first die described above is thrown twice and the two numbers on the die noted. Is\par
it possible to find values of q so that the probability that the numbers are the same is\par
j\par
\[less than 1/36?\]
13 Bar magnets are placed randomly end-to-end in a straight line. If adjacent magnets have\par
ends of opposite polarities facing each other, they join together to form a single unit. If they\par
have ends of the same polarity facing each other, they stand apart. Find the expectation and\par
variance of the number of separate units in terms of the total number N of magnets.\par
14 When I throw a dart at a target, the probability that it lands a distance X from the centre is\par
a random variable with density function\par
2x if 0 6 x 6 1 ;\par
\[f (x) =\]
( 0 otherwise.\par
I score points according to the position of the dart as follows: if 0 6 X < 1 , my score is 4;\par
\[if 1 6 X < 1 , my score is 3; if 1 6 X < 3 , my score is 2; if 3 6 X 6 1, my score is 1.\]
4 2 2 4 4\par
\textbf{(i)} Show that my expected score from one dart is 15/8.\par
\textbf{(ii)} I play a game with the following rules. I start off with a total score 0, and each time I\par
throw a dart my score on that throw is added to my total. Then:\par
if my new total is greater than 3, I have lost and the game ends;\par
if my new total is 3, I have won and the game ends;\par
if my new total is less than 3, I throw again.\par
Show that, if I have won such a game, the probability that I threw the dart three\par
\[times is 343/2231.\]
\end{flushleft}
\end{document}
